两个聚合度相等 \(N\ge1\) 的聚合物A、B在不可压格点模型中混合，A的体积分数 \(0<\phi<1\)。无量纲每格点Flory–Huggins混合自由能为
\[
g(\phi)=\frac{\phi}{N}\ln\phi+\frac{1-\phi}{N}\ln(1-\phi)+\chi\phi(1-\phi),
\]
其中相互作用参数 \(\chi\) 可随温度改变，忽略梯度能和动力学。
(1) 求混合自由能的二阶导数与局部稳定性，得到旋节线 \(\chi_s(\phi)\)。
(2) 用临界点二阶、三阶导数同时为零求临界浓度与临界 \(\chi_c\)，检验四阶导数的稳定号。
(3) 利用对称性给出两相共存的binodal隐式方程，并在临界附近推导组分差随 \(\chi-\chi_c\) 的平方根标度。
(4) 比较binodal与spinodal的含义，讨论 \(\chi\to0\)、\(N\to\infty\) 和有限尺寸时相分离判据的适用范围。
